\section*{अध्याय \ref{ch:g10:the-mole} --- मोल और मोलर द्रव्यमान}

\begin{solution}{exo:g10:the-mole:1}
$M(\ce{O2}) = 2 \times 16.0 = \qty{32.0}{g/mol}$;
$M(\ce{N2}) = 2 \times 14.0 = \qty{28.0}{g/mol}$;
$M(\ce{CH4}) = 12.0 + 4 \times 1.0 = \qty{16.0}{g/mol}$;
$M(\ce{NH3}) = 14.0 + 3 \times 1.0 = \qty{17.0}{g/mol}$;
$M(\ce{CO2}) = 12.0 + 2 \times 16.0 = \qty{44.0}{g/mol}$।
\end{solution}

\begin{solution}{exo:g10:the-mole:2}
पानी: $n = 36.0 / 18.0 = \qty{2.00}{mol}$। मेथेन:
$n = 4.0 / 16.0 = \qty{0.25}{mol}$।
\end{solution}

\begin{solution}{exo:g10:the-mole:3}
$m(\ce{NaCl}) = 0.25 \times 58.5 = \qty{14.6}{g}$, लगभग \qty{15}{g};
$m(\ce{CO2}) = 2.0 \times 44.0 = \qty{88}{g}$।
\end{solution}

\begin{solution}{exo:g10:the-mole:4}
$N = 0.50 \times \num{6.02e23} = \num{3.0e23}$ डाइऑक्सीजन \omterm{def:g7:atoms-and-molecules:molecule}{अणु}।
$n = \num{3.0e24} / \qty{6.02e23}{mol^{-1}} = \qty{5.0}{mol}$ लोहा।
\end{solution}

\begin{solution}{exo:g10:the-mole:5}
$V = n \times V_m = 0.50 \times 24.1 = \qty{12}{L}$।
$n = V / V_m = 12 / 24.1 = \qty{0.50}{mol}$।
\end{solution}

\begin{solution}{exo:g10:the-mole:6}
बूँद का द्रव्यमान \qty{0.050}{g} है, इसलिए
$n = 0.050 / 18.0 = \qty{2.8e-3}{mol}$ और
$N = \num{2.8e-3} \times \num{6.02e23} = \num{1.7e21}$ \omterm{def:g7:atoms-and-molecules:molecule}{अणु}।
\end{solution}

\begin{solution}{exo:g10:the-mole:7}
ऐलुमिनियम \omterm{def:g7:atoms-and-molecules:atom}{परमाणु} (\qty{27.0}{g/mol}) लोहे के \omterm{def:g7:atoms-and-molecules:atom}{परमाणु} (\qty{55.8}{g/mol}) से
लगभग आधा भारी है, इसलिए ऐलुमिनियम के उतने ही द्रव्यमान में लगभग
दुगने \omterm{def:g7:atoms-and-molecules:atom}{परमाणु} होते हैं। जाँच: ऐलुमिनियम $10/27.0 = \qty{0.370}{mol}$, यानी
\num{2.23e23} \omterm{def:g7:atoms-and-molecules:atom}{परमाणु}; लोहा $10/55.8 = \qty{0.179}{mol}$, यानी
\num{1.08e23} \omterm{def:g7:atoms-and-molecules:atom}{परमाणु}। ऐलुमिनियम आगे है, लगभग 2.1 गुना।
\end{solution}

\begin{solution}{exo:g10:the-mole:8}
एक लीटर गैस $1/24.1 = \qty{0.0415}{mol}$ है। कार्बन डाइऑक्साइड:
$0.0415 \times 44.0 = \qty{1.83}{g}$; डाइनाइट्रोजन:
$0.0415 \times 28.0 = \qty{1.16}{g}$। कार्बन डाइऑक्साइड \omterm{def:g4:air-a-mixture-of-gases:air}{वायु} से लगभग डेढ़ गुना
घनी है, जिसका \omterm{def:g10:the-mole:molar-mass}{मोलर द्रव्यमान} डाइनाइट्रोजन के क़रीब है। किण्वित होता रस
कार्बन डाइऑक्साइड छोड़ता है, जो नीचे बैठकर बंद तहख़ाने के फ़र्श के पास जमा होती
जाती है, जहाँ वह अंदर जाने वाले किसी का भी दम घोंट सकती है।
\end{solution}

\begin{solution}{exo:g10:the-mole:9}
$M(\ce{C6H12O6}) = 6 \times 12.0 + 12 \times 1.0 + 6 \times 16.0 =
\qty{180.0}{g/mol}$। तीर $\div N_A$: $n = \num{1.0e22} / \num{6.02e23} =
\qty{1.66e-2}{mol}$; तीर $\times M$: $m = \num{1.66e-2} \times 180.0 =
\qty{3.0}{g}$।
\end{solution}

\begin{solution}{exo:g10:the-mole:10}
$M = \qty{342.0}{g/mol}$, इसलिए $n = 5.0 / 342.0 = \qty{1.46e-2}{mol}$;
$N = \num{1.46e-2} \times \num{6.02e23} = \num{8.8e21}$ \omterm{def:g7:atoms-and-molecules:molecule}{अणु}। हर एक में
12 कार्बन \omterm{def:g7:atoms-and-molecules:atom}{परमाणु} हैं: $12 \times \num{8.8e21} = \num{1.1e23}$ कार्बन \omterm{def:g7:atoms-and-molecules:atom}{परमाणु}।
\end{solution}

\begin{solution}{exo:g10:the-mole:11}
पानी: $1000 / 18.0 = \qty{55.6}{mol}$। एथेनॉल:
$M = 2 \times 12.0 + 6 \times 1.0 + 16.0 = \qty{46.0}{g/mol}$, इसलिए
$1000 / 46.0 = \qty{21.7}{mol}$। पानी की मात्रा अधिक है,
$55.6 / 21.7 \approx 2.6$ गुना, जो बस $46.0 / 18.0$ है।
\end{solution}

\begin{solution}{exo:g10:the-mole:12}
सोना: $\qty{3.75}{g}$, $n = 3.75 / 197.0 = \qty{1.90e-2}{mol}$, यानी
\num{1.15e22} \omterm{def:g7:atoms-and-molecules:atom}{परमाणु}। ताँबा: \qty{1.25}{g},
$n = 1.25 / 63.5 = \qty{1.97e-2}{mol}$, यानी \num{1.19e22} \omterm{def:g7:atoms-and-molecules:atom}{परमाणु}।
आश्चर्य की बात है कि अँगूठी में सोने के \omterm{def:g7:atoms-and-molecules:atom}{परमाणुओं} से थोड़े अधिक ताँबे के \omterm{def:g7:atoms-and-molecules:atom}{परमाणु} हैं:
सोने का \omterm{def:g7:atoms-and-molecules:atom}{परमाणु} ताँबे के \omterm{def:g7:atoms-and-molecules:atom}{परमाणु} से लगभग तीन गुना भारी है।
\end{solution}

\begin{solution}{exo:g10:the-mole:13}
\begin{enumerate}
  \item $V = 8.0 \times 6.0 \times 3.0 = \qty{144}{m^3} =
        \qty{1.44e5}{L}$; $n = \num{1.44e5} / 24.1 = \qty{5.98e3}{mol}$;
        $N = \num{5.98e3} \times \num{6.02e23} = \num{3.6e27}$ \omterm{def:g7:atoms-and-molecules:molecule}{अणु}।
  \item $M = (78 \times 28.0 + 21 \times 32.0 + 1 \times 40.0) / 100 =
        2896 / 100 = \qty{29.0}{g/mol}$।
  \item $m = \num{5.98e3} \times 29.0 = \qty{1.73e5}{g}$, लगभग
        \qty{173}{kg}: एक कक्षा की \omterm{def:g4:air-a-mixture-of-gases:air}{वायु} का भार दो-तीन वयस्कों
        जितना है।
\end{enumerate}
\end{solution}

\begin{solution}{exo:g10:the-mole:14}
\begin{enumerate}
  \item $n = 1000 / 28.1 = \qty{35.6}{mol}$, इसलिए
        $N = 35.6 \times \num{6.02e23} = \num{2.14e25}$ \omterm{def:g7:atoms-and-molecules:atom}{परमाणु}।
  \item $m = \qty{1.000}{kg} / \num{2.14e25} = \qty{4.67e-26}{kg}$।
  \item गोले का द्रव्यमान और \omterm{def:g10:the-mole:molar-mass}{मोलर द्रव्यमान} उसकी मात्रा देते हैं,
        $n = m / M$। अगर \omterm{def:g7:atoms-and-molecules:atom}{परमाणुओं} $N$ को स्वतंत्र रूप से गिना जाए (गोले के
        आयतन और क्रिस्टल में \omterm{def:g7:atoms-and-molecules:atom}{परमाणुओं} के बीच की दूरी से), तो अनुपात $N / n$
        \omterm{def:g10:the-mole:avogadro}{आवोगाद्रो स्थिरांक} है।
\end{enumerate}
\end{solution}

\begin{solution}{exo:g10:the-mole:15}
$M = 0.758 \times 35.0 + 0.242 \times 37.0 = 26.53 + 8.95 =
\qty{35.5}{g/mol}$: सारणी वाला मान।
\end{solution}

\begin{solution}{pb:g10:the-mole:1}
\textbf{1.} $M(\ce{H2O}) = 2 \times 1.0 + 16.0 = \qty{18.0}{g/mol}$।

\textbf{2.} $m = 250 \times \qty{1.0}{g} = \qty{250}{g}$।

\textbf{3.} $n = 250 / 18.0 = \qty{13.9}{mol}$।

\textbf{4.} हर \omterm{def:g7:atoms-and-molecules:molecule}{अणु} में दो हाइड्रोजन \omterm{def:g7:atoms-and-molecules:atom}{परमाणु} और एक ऑक्सीजन \omterm{def:g7:atoms-and-molecules:atom}{परमाणु} हैं:
हाइड्रोजन \omterm{def:g7:atoms-and-molecules:atom}{परमाणुओं} का \qty{27.8}{mol}, ऑक्सीजन \omterm{def:g7:atoms-and-molecules:atom}{परमाणुओं} का \qty{13.9}{mol}।

\textbf{5.} $N = 13.9 \times \num{6.02e23} = \num{8.36e24}$ \omterm{def:g7:atoms-and-molecules:molecule}{अणु}।

\textbf{6.} $m = \qty{18.0}{g/mol} / \qty{6.02e23}{mol^{-1}} =
\qty{2.99e-23}{g}$।

\textbf{7.} $\num{8.36e24} \times \qty{2.99e-23}{g} = \qty{250}{g}$, जो
गिलास के पानी का द्रव्यमान है, जैसा होना चाहिए।

\textbf{8.} $\num{8.36e24} / \num{3.16e7} = \num{2.6e17}$ वर्ष: सत्रह शून्यों
वाली वर्षों की संख्या, किसी भी संभव गिनती से बहुत परे।

\textbf{9.} $\qty{1.335e9}{km^3} \times \num{e9} = \qty{1.335e18}{m^3}$,
और $\qty{1}{m^3} = \qty{1000}{L}$ से, \qty{1.335e21}{L}।

\textbf{10.} $\num{8.36e24} / \num{1.335e21} = \num{6.26e3}$ निशान वाले
\omterm{def:g7:atoms-and-molecules:molecule}{अणु} प्रति लीटर।

\textbf{11.} $\num{6.26e3} \times 0.250 = \num{1.57e3}$ निशान वाले \omterm{def:g7:atoms-and-molecules:molecule}{अणु}
दूसरे गिलास में।

\textbf{12.} द्रव्यमान $\num{1.335e21} \times \qty{1.0}{kg} =
\qty{1.335e21}{kg} = \qty{1.335e24}{g}$; $n = \num{1.335e24} / 18.0 =
\qty{7.42e22}{mol}$।

\textbf{13.} $13.9 / \num{7.42e22} = \num{1.87e-22}$।

\textbf{14.} $\num{1.87e-22} \times \num{8.36e24} = \num{1.57e3}$: प्रश्न 11
वाला ही उत्तर, जैसा होना चाहिए।

\textbf{15.} $\num{1.335e21} / 0.250 = \num{5.34e21}$ गिलास।

\textbf{16.} $\num{8.36e24} / \num{5.34e21} \approx 1570$। एक गिलास में
जितने गिलास पानी महासागरों में है, उससे लगभग 1600 गुना अधिक \omterm{def:g7:atoms-and-molecules:molecule}{अणु} हैं;
समान रूप से फैलने पर एक गिलास के \omterm{def:g7:atoms-and-molecules:molecule}{अणु} महासागर के हर गिलास भर पानी में
लगभग 1600 रह जाते हैं।

\textbf{17.} $1 / \num{6.26e3} = \qty{1.6e-4}{L} = \qty{0.16}{mL}$, लगभग
तीन बूँदें।

\textbf{18.} महासागर का आयतन एक अनुमान है, जो अधिक से अधिक तीन-चार अंकों तक ज्ञात
है, और उसे पूर्णांकित किया गया; समुद्र का पानी शुद्ध पानी नहीं है (उसमें
लवण हैं); गिलास ठीक \qty{250}{mL} का नहीं है; और सारे महासागरों में मिलना
कभी पूरी तरह समान नहीं होगा। केवल परिमाण की कोटि पर भरोसा किया
जा सकता है।

\textbf{19.} पहले गिलास के लगभग 1600 \omterm{def:g7:atoms-and-molecules:molecule}{अणु} दूसरे में वापस
आते हैं।
\end{solution}
